\section{Regularity and sharpness of the dimension bound}
\label{sec:consequences}

We now pass from the classification of homogeneous global minimizers to
local free boundaries. This step uses the dimension-reduction theorem of
Weiss \cite{Weiss1999}, in the formulation
\cite[Theorem~1.4 and Proposition~10.13]{Velichkov2023}.

\begin{proof}[Proof of Corollary~\ref{cor:singular-set}]
Apply the cited dimension-reduction theorem with volume coefficient one
on balls compactly contained in $D$. Its inputs are uniform flatness
regularity, linear nondegeneracy, compactness of blowups, and the fact that
all blowup limits are homogeneous global minimizers. These are the
constant-coefficient results used in Proposition~\ref{red:regularity};
flatness regularity is \cite[Theorem~8.1]{Velichkov2023}.
The theorem gives an empty singular set below $d_*$, local finiteness in
dimension $d_*$, and Hausdorff dimension at most $n-d_*$ above it.
Theorem~\ref{thm:main} sets $d_*=7$. Smoothness at each regular point
follows from Proposition~\ref{red:regularity}(iii). A countable cover by
interior balls gives the assertions on $D$.

For sharpness, let $U$ be the nonflat global minimizer in $\R^7$ of
De Silva and Jerison \cite[Theorem~1.1]{DeSilvaJerison2009}, and set
$U_n(x,y)=U(x)$ on $\R^7\times\R^{n-7}$.
Repeated application of Proposition~\ref{red:suspension} shows that $U_n$
is globally minimizing. At each point $(0,y_0)$, every rescaling of $U_n$
is $U_n$ itself, by homogeneity and translation invariance in $y$.
It cannot be flat: a half-space normal would have zero $y$-component,
forcing $U$ to be flat. Thus
$\{0\}\times\R^{n-7}\subset\operatorname{Sing}(U_n)$.
Its dimension is $n-7$, and the proved upper bound gives equality.
\end{proof}
